% Einheit 2 von 8 – Laplace-Experimente in der Oberstufe (bank/zufallsexperimente-und-pfadregeln/e2.jsonl, zusammenbau v0.1)
%% TODO Zeitmarke Einheit 2: Marken-Zeile nicht lesbar
%% TODO Zweigzeile Teil 1: was hier gelernt wird (Satz in Schülersprache aus dem Einheitstitel) – Einheit 2
\einheitenkopf[e2]{Einheit 2 von 8 · Laplace-Experimente in der Oberstufe}
\zweigzeile{Abitur GK}
\setcounter{aufgabe}{13}

%% TODO Titel: Ich-kann-Satz fehlt in der Bank (Platzhalter: Kettenname) – Nr. 14
%% TODO Anweisung: ein Satz über der ersten Teilaufgabe fehlt in der Bank – Nr. 14
\begin{aufgabe}{Laplace-Experimente}
%% TODO \rechenplatz{n}: Schrittzahl des Grundfalls fehlt in der Bank (nur bei fester Schreibform, 2.3 b)
\begin{teile}
\teil Ein Glücksrad hat vier gleich große Sektoren. Sind die vier Ergebnisse gleich wahrscheinlich?\\ \kreuz{ja, Laplace}\\ \kreuz{nein, kein Laplace}
\teil In einer Lostrommel liegen 50 Lose, 8 davon sind Gewinne. P(Gewinn)? P = \leerfeld
\teil Aus einem Beutel mit 30 gleich großen, gleich schweren Murmeln wird blind eine gezogen. Begründe, dass das ein Laplace-Experiment ist, und nenne ein weiteres Beispiel.
\teil Eine Münze wird dreimal geworfen. P(genau zweimal Kopf)? P = \leerfeld
\teil Ein roter und ein blauer Würfel werden geworfen. Weise mit einer Ergebnistabelle nach: P(Augensumme 9) $= \frac{1}{9}$.
\teil Zwei Würfel werden geworfen. Augensumme 5 oder Augensumme 9 – was ist wahrscheinlicher? Begründe ohne Rechnung.
\teil Ein Glücksrad besteht aus gleich großen Sektoren von je $10^\circ$. P(Gewinn) soll $\frac{1}{4}$ sein, P(Freilos) $\frac{1}{9}$, der Rest ist Niete. Wie viele Sektoren je Art? \leerfeld
\teil Behälter X enthält 4 rote Kugeln, Behälter Y 2 rote und 5 blaue. Eine Kugel wird zufällig von Y nach X gelegt. P(danach sind in X so viele rote wie in Y blaue)? Begründe.
\teil Eine Urne enthält n Kugeln, 20\,\% davon tragen einen Stern. Zwei Kugeln mit Stern werden hinzugelegt, dann wird eine Kugel gezogen. Weise über die Differenz der Wahrscheinlichkeiten nach, dass P(Stern) größer geworden ist. \hfill (Abitur 2019 GK)
\end{teile}
\end{aufgabe}

%% TODO Titel: Ich-kann-Satz fehlt in der Bank (Platzhalter: Kettenname) – Nr. 15
%% TODO Anweisung: ein Satz über der ersten Teilaufgabe fehlt in der Bank – Nr. 15
\begin{aufgabe}{Laplace-Experiment: Aussage über Wahrscheinlichkeiten aus absoluten Anzahlen ohne Grundgesamtheit beurteilen}
\begin{teile}
\teil In Stadt A wurden im letzten Jahr 1\,200 Fahrräder gestohlen, in Stadt B 600. Aussage: „In A ist die Wahrscheinlichkeit, dass ein Rad gestohlen wird, doppelt so groß.“ Beurteile.
\end{teile}
\end{aufgabe}

%% TODO Titel: Ich-kann-Satz fehlt in der Bank (Platzhalter: Kettenname) – Nr. 16
%% TODO Anweisung: ein Satz über der ersten Teilaufgabe fehlt in der Bank – Nr. 16
\begin{aufgabe}{Laplace-Experimente – Fehler finden, Begründen, Anwendung, Darstellungswechsel}
\begin{teile}
\teil Zwei Würfel: P(erster Wurf kleiner als zweiter)? Jan zählt 21 günstige Paare und schreibt $P = \frac{21}{36}$. Finde den Fehler.
\teil Begründe: Beim Vergleich zweier Laplace-Wahrscheinlichkeiten desselben Experiments darf man den Nenner weglassen.
\teil Für das Schulfest wird ein Glücksrad aus 24 gleich großen Sektoren gebaut. Der Hauptgewinn soll mit $\frac{1}{12}$, ein Trostpreis mit $\frac{1}{3}$ fallen, sonst gibt es nichts. Wie viele Sektoren je Art? \leerfeld
\teil Ein Glücksrad hat Sektoren mit $90^\circ$ rot, $90^\circ$ blau und $180^\circ$ gelb. Beschreibe eine Urne mit möglichst wenigen Kugeln, die dasselbe leistet. \leerfeld
\end{teile}
\end{aufgabe}
